\documentclass[11pt,letterpaper]{article}

\usepackage[spanish,es-noshorthands]{babel}
\usepackage{fontspec}
\setmainfont{Source Serif Pro}
\usepackage{microtype}
\usepackage{geometry}
\usepackage{booktabs}
\usepackage{tabularx}
\usepackage{enumitem}
\usepackage{xcolor}
\usepackage{graphicx}
\usepackage{csquotes}
\usepackage{amsmath}
\usepackage{siunitx}
\usepackage[version=4]{mhchem}
\usepackage{tikz}
\usetikzlibrary{arrows.meta,positioning,shapes.geometric,fit,backgrounds,patterns}
\usepackage{xurl}
\usepackage[hidelinks]{hyperref}
\usepackage[backend=biber,style=apa,sorting=nyt]{biblatex}
\DeclareLanguageMapping{spanish}{spanish-apa}

\geometry{margin=2.3cm}
% Los rótulos de las figuras usan nodos de ancho fijo con texto alineado, y la
% bibliografía compone en modo sloppy: ambos generan avisos de línea holgada
% que son cosméticos. Se silencian para que `make check` siga siendo sensible
% a los desbordes reales.
\hbadness=10000
\setlength{\emergencystretch}{1em}
\setlength{\parindent}{0pt}
\setlength{\parskip}{0.55em}
\setlist{nosep}
\renewcommand{\topfraction}{0.92}
\renewcommand{\bottomfraction}{0.85}
\renewcommand{\textfraction}{0.06}
\renewcommand{\floatpagefraction}{0.72}
\setcounter{topnumber}{2}
\setcounter{totalnumber}{3}
\addbibresource{referencias.bib}
\AtBeginBibliography{\sloppy}

\sisetup{
  output-decimal-marker = {,},
  group-separator = {.},
  group-minimum-digits = 4,
  detect-all
}
\DeclareSIUnit{\pkm}{pax\text{-}km}
\DeclareSIUnit{\pax}{pax}

% Colores: subsónico, supersónico, tiempo en tierra, dato de referencia y
% señal de alerta o límite
\definecolor{cSub}{RGB}{16,110,104}
\definecolor{cSup}{RGB}{186,98,52}
\definecolor{cTierra}{RGB}{132,124,112}
\definecolor{cRef}{RGB}{70,90,112}
\definecolor{cAviso}{RGB}{176,42,42}
\definecolor{cVerde}{RGB}{74,116,60}

% --- Figura 1: velocidad de crucero y consumo -------------------------------
\newcommand{\vAx}[1]{(#1-1934)*0.150}
\newcommand{\vAy}[1]{#1*0.00250}
% #1 año, #2 velocidad, #3 color, #4 etiqueta, #5 posición del nodo
\newcommand{\avion}[5]{%
  \fill[#3] ({\vAx{#1}},{\vAy{#2}}) circle (2.1pt);%
  \node[#5, font=\scriptsize, inner sep=2.6pt, text=#3]
       at ({\vAx{#1}},{\vAy{#2}}) {#4};%
}
\newcommand{\vBx}[1]{(#1-1956)*0.190}
\newcommand{\vBy}[1]{#1*0.0320}

% --- Figura 3: eficiencia frente a tráfico ----------------------------------
\newcommand{\eAx}[1]{(#1-2009)*0.385}
\newcommand{\eAy}[1]{#1*0.0160}

% --- Figura 2: barras apiladas de tiempo puerta a puerta --------------------
% #1 y, #2 h en tierra antes, #3 h de vuelo, #4 h en tierra después,
% #5 rótulo, #6 color del tramo de vuelo, #7 total puerta a puerta
\newcommand{\barrapp}[7]{%
  \node[left, text width=3.1cm, align=right, font=\scriptsize]
       at (-0.12,#1) {#5};%
  \fill[cTierra!45] (0,#1-0.26) rectangle ({#2*1.02},#1+0.26);%
  \fill[#6] ({#2*1.02},#1-0.26) rectangle ({(#2+#3)*1.02},#1+0.26);%
  \fill[cTierra!45] ({(#2+#3)*1.02},#1-0.26)
                    rectangle ({(#2+#3+#4)*1.02},#1+0.26);%
  \node[font=\scriptsize\bfseries, text=white]
       at ({(#2+#3/2)*1.02},#1) {\num{#3} h};%
  \node[right, font=\scriptsize] at ({(#2+#3+#4)*1.02+0.14},#1)
       {\textbf{#7 h}};%
}

% --- Figura 3: barras verticales sencillas ----------------------------------
% #1 x, #2 alto, #3 color, #4 rótulo bajo el eje, #5 valor sobre la barra
\newcommand{\barrav}[5]{%
  \fill[#3] (#1-0.42,0) rectangle (#1+0.42,#2);%
  \node[above, font=\scriptsize\bfseries, text=#3] at (#1,#2) {#5};%
  \node[below, font=\scriptsize, text width=2.9cm, align=center]
       at (#1,-0.08) {#4};%
}

% --- Figura 4: barras horizontales normalizadas -----------------------------
% #1 y, #2 fin de la barra, #3 color, #4 rótulo izquierdo, #5 valor
\newcommand{\barrah}[5]{%
  \node[left, font=\scriptsize] at (-0.12,#1) {#4};%
  \fill[#3] (0,#1-0.155) rectangle (#2,#1+0.155);%
  \node[right, font=\scriptsize] at ({#2+0.1},#1) {#5};%
}

\title{\vspace{-1.2cm}\textbf{Mi experiencia volando, leída con las figuras de
mérito del módulo}\\
\large Seis respuestas personales y lo que la evolución tecnológica del
transporte aéreo dice de cada una}
\author{Carlos Luis Rojas Aragonés\\
\small Gestión del Desarrollo Tecnológico}
\date{\small 7 de septiembre de 2026}

\begin{document}

\maketitle
\vspace{-0.6em}

\section{Cómo respondo a estas seis preguntas}

Las seis preguntas del foro son personales y las contesto como tales, sin
adornarlas. Pero el módulo acaba de darnos un instrumento que sirve
exactamente para esto: la \textbf{figura de mérito}, es decir, la medida
cuantitativa del desempeño de una tecnología en la función que presta,
independiente de la implementación concreta que la entregue
\parencite{deweck2022}. Cada una de mis respuestas contiene, escondida, una
figura de mérito que estoy usando sin declararla. Este trabajo consiste en
declararla y en comprobar qué dice la evidencia sobre ella.

El resultado más útil de hacerlo es que \textbf{en tres de las seis respuestas
la figura de mérito que yo tenía en la cabeza no era la que de verdad me
importa}. En la de la velocidad, sobre todo. Ese desajuste no es un error mío:
es el mismo desajuste que explica por qué el avión en el que vuelo hoy va a la
misma velocidad que el que volaba mi abuelo.

El Cuadro~\ref{cuadro:mapa} adelanta el mapa completo.

\begin{table}[t]
\small
\caption{Las seis preguntas, mi respuesta y la figura de mérito implícita.}
\label{cuadro:mapa}
\begin{tabularx}{\textwidth}{@{}p{3.5cm}p{4.1cm}X@{}}
\toprule
\textbf{Pregunta} & \textbf{Mi respuesta} & \textbf{Figura de mérito que estoy
usando en realidad} \\
\midrule
Cuántas veces vuelo al año
& Seis, todas por trabajo
& Ninguna: es el dato de entrada que hace medibles las otras cinco \\
Cómo mejoraría la experiencia
& Que el vuelo dure la mitad
& Creí que era la velocidad de crucero. Es el \emph{tiempo sentado}, que no es
lo mismo \\
Vergüenza de volar
& No me afecta; mis viajes son necesarios
& \si{\kilogram} de \ce{CO2} equivalente por pasajero-kilómetro, y su tasa de
mejora frente al crecimiento del tráfico \\
Los tres mejores aeropuertos
& Atlanta y Panamá
& Dos distintas y las dos legítimas: tránsito y conectividad de red \\
Los tres que más frecuento
& \mbox{Juan Santamaría}, Houston y Los Ángeles
& Ninguna de calidad: es la topología de la red vista desde San José \\
Modos más lentos y limpios
& No me preocupa; que avance la tecnología
& Tiempo puerta a puerta frente a \ce{CO2} por viaje, en un frente de Pareto
que en Costa Rica tiene un solo punto \\
\bottomrule
\end{tabularx}
\end{table}

\section{Cuántas veces vuelo al año}

\textbf{Este año he volado seis veces, todas por motivos de trabajo.} Ninguna
de esas seis fue por placer, y ninguna la habría hecho si el asunto se hubiera
podido resolver por videoconferencia.

Seis es un número modesto en términos absolutos y muy alto en términos
globales: alrededor del 80 por ciento de la población mundial no ha volado
nunca \parencite{owid2026aviacion}. Lo anoto no como confesión sino porque el
número importa para las preguntas 3 y 6, que solo se pueden contestar con
honestidad si uno pone antes su propia cifra sobre la mesa.

Como no llevo registro de itinerarios, en el Cuadro~\ref{cuadro:rutas} calculo
las distancias de círculo máximo de las rutas que sí uso y las convierto en
tiempo y en emisiones. Las distancias son cálculo propio a partir de las
coordenadas de cada aeropuerto; los factores de emisión son los oficiales
británicos, que son los que usa la mayoría de las herramientas de contabilidad
corporativa \parencite{defra2025}.

\begin{table}[t]
\small
\caption{Mis rutas habituales: distancia, tiempo de bloque y emisiones por
tramo y pasajero en clase turista. Distancias de círculo máximo, cálculo
propio. Tiempo de bloque modelado como \SI{0.6}{\hour} fijas de rodaje,
ascenso y descenso más la distancia dividida entre \SI{903}{\kilo\meter\per\hour}
(Mach 0,85 a altitud de crucero). Factores de emisión de
\textcite{defra2025}: \SI{0.150}{\kilogram} por pasajero-kilómetro hasta
\SI{3700}{\kilo\meter} y \SI{0.117}{\kilogram} por encima.}
\label{cuadro:rutas}
\begin{tabularx}{\textwidth}{@{}Xrrrr@{}}
\toprule
\textbf{Ruta} & \textbf{Distancia} & \textbf{Tiempo de} &
\textbf{Factor} & \textbf{\ce{CO2}e por} \\
 & \textbf{(\si{\kilo\meter})} & \textbf{bloque (\si{\hour})} &
 \textbf{(\si{\kilogram\per\pkm})} & \textbf{tramo (\si{\kilogram})} \\
\midrule
San José--Panamá (SJO--PTY)          &  539 & 1,2 & 0,150 &  81 \\
San José--Houston (SJO--IAH)         & 2.505 & 3,4 & 0,150 & 376 \\
San José--Atlanta (SJO--ATL)         & 2.629 & 3,5 & 0,150 & 394 \\
San José--Los Ángeles (SJO--LAX)     & 4.382 & 5,5 & 0,117 & 513 \\
\bottomrule
\end{tabularx}
\end{table}

Seis tramos de esas rutas suman entre \textbf{2,3 y 3,1 toneladas} de \ce{CO2}
equivalente, según la mezcla: 2,3 si los seis fueron a Houston, 3,1 si los seis
fueron a Los Ángeles. Volveré sobre esa cifra en la pregunta 3, porque es la
única forma seria de contestarla.

\section{Cómo creo que podría mejorar mi experiencia}
\label{sec:velocidad}

Mi respuesta espontánea fue esta: \textbf{que los vuelos fueran mucho más
rápidos, que duraran la mitad o menos.} Me desespero después de más de cinco
horas metido en un avión.

Es la respuesta que da casi todo el mundo, y resulta que es la más interesante
de las seis, porque contiene dos cosas que conviene separar: una figura de
mérito que lleva sesenta y ocho años sin moverse y un diagnóstico equivocado de
mi propia molestia.

\subsection{La figura de mérito que dejó de mejorar}

La velocidad de crucero del avión comercial es el ejemplo de manual de una
figura de mérito que se satura. La Figura~\ref{fig:velocidad} la dibuja.

\begin{figure}[t]
\centering
\begin{tikzpicture}

% ---------- Panel A: velocidad de crucero ----------
\node[anchor=south west, font=\small\bfseries] at (0,6.15)
     {Panel A. Velocidad de crucero del avión comercial, 1936--2030};

% rejilla y ejes
\foreach \v in {500,1000,1500,2000} {
  \draw[cTierra!30] (0,{\vAy{\v}}) -- (15.0,{\vAy{\v}});
  \node[left, font=\scriptsize, text=cTierra] at (-0.06,{\vAy{\v}}) {\v};
}
\draw[cTierra!70] (0,0) -- (15.0,0);
\draw[cTierra!70] (0,0) -- (0,5.85);
\foreach \a in {1940,1950,1960,1970,1980,1990,2000,2010,2020,2030} {
  \draw[cTierra!70] ({\vAx{\a}},0) -- ({\vAx{\a}},-0.09);
  \node[below, font=\scriptsize, text=cTierra] at ({\vAx{\a}},-0.09) {\a};
}
\node[left, font=\scriptsize, text=cTierra, rotate=90, anchor=south]
     at (-0.75,2.6) {km/h};

% banda de la meseta subsónica
\fill[cSub!10] ({\vAx{1958}},{\vAy{870}})
     rectangle ({\vAx{2032}},{\vAy{925}});
\node[right, font=\scriptsize\itshape, text=cSub!85!black]
     at ({\vAx{2004}},{\vAy{800}}) {meseta: 68 años, \SI[retain-explicit-plus]{+0.7}{\percent}};

% trayectoria subsónica
\draw[cSub, thick]
  ({\vAx{1936}},{\vAy{333}}) -- ({\vAx{1945}},{\vAy{480}}) --
  ({\vAx{1947}},{\vAy{507}}) -- ({\vAx{1956}},{\vAy{555}}) --
  ({\vAx{1958}},{\vAy{897}}) -- ({\vAx{1970}},{\vAy{907}}) --
  ({\vAx{1995}},{\vAy{892}}) -- ({\vAx{2007}},{\vAy{903}}) --
  ({\vAx{2011}},{\vAy{903}}) -- ({\vAx{2018}},{\vAy{903}});

\avion{1936}{333}{cSub}{DC-3}{below right}
\avion{1945}{480}{cSub}{Constellation}{above left}
\avion{1947}{507}{cSub}{DC-6}{below right}
\avion{1952}{740}{cSub!60!cRef}{Comet 1}{above left}
\avion{1956}{555}{cSub}{DC-7C}{below right}
\avion{1958}{897}{cSub}{707}{above left}
\avion{1970}{907}{cSub}{747}{above}
\avion{1995}{892}{cSub}{777}{above}
\avion{2011}{903}{cSub}{787}{above}
\avion{2018}{903}{cSub}{A350}{above right}

% Concorde y su retirada
\avion{1976}{2179}{cSup}{Concorde, Mach 2,02}{above}
\draw[cSup, dashed, thick, -{Stealth[length=4pt]}]
  ({\vAx{1976}},{\vAy{2120}}) -- ({\vAx{2003}},{\vAy{2120}})
  -- ({\vAx{2003}},{\vAy{960}});
\node[right, font=\scriptsize, text=cSup]
     at ({\vAx{1978}},{\vAy{1930}}) {27 años en servicio; retirado en 2003};

% Overture proyectado
\draw[cSup, thick, dotted] ({\vAx{2018}},{\vAy{903}})
  -- ({\vAx{2030}},{\vAy{1806}});
\draw[cSup, fill=white, thick] ({\vAx{2030}},{\vAy{1806}}) circle (2.4pt);
\node[left, font=\scriptsize, text=cSup, align=right]
     at ({\vAx{2029}},{\vAy{1480}}) {Overture, Mach 1,7\\(previsto para 2030)};

% ---------- Panel B: energía por pasajero-kilómetro ----------
\begin{scope}[yshift=-5.7cm]
\node[anchor=south west, font=\small\bfseries] at (0,3.75)
     {Panel B. La figura de mérito que sí siguió mejorando: energía por
      pasajero-kilómetro};

\foreach \v in {25,50,75,100} {
  \draw[cTierra!30] (0,{\vBy{\v}}) -- (14.1,{\vBy{\v}});
  \node[left, font=\scriptsize, text=cTierra] at (-0.06,{\vBy{\v}}) {\v};
}
\draw[cTierra!70] (0,0) -- (14.1,0);
\draw[cTierra!70] (0,0) -- (0,3.55);
\foreach \a in {1960,1970,1980,1990,2000,2010,2020,2030} {
  \draw[cTierra!70] ({\vBx{\a}},0) -- ({\vBx{\a}},-0.09);
  \node[below, font=\scriptsize, text=cTierra] at ({\vBx{\a}},-0.09) {\a};
}
\node[left, font=\scriptsize, text=cTierra, rotate=90, anchor=south]
     at (-0.8,1.7) {índice, 1958 = 100};

\draw[cVerde, very thick, samples=80, domain=1958:2026, variable=\x]
  plot ({\vBx{\x}},{(100*pow(0.98,\x-1958))*0.0320});
\draw[cVerde, thick, dashed, samples=20, domain=2026:2030, variable=\x]
  plot ({\vBx{\x}},{(100*pow(0.98,\x-1958))*0.0320});

\fill[cVerde] ({\vBx{1958}},{\vBy{100}}) circle (2.1pt);
\node[above right, font=\scriptsize, text=cVerde]
     at ({\vBx{1958}},{\vBy{100}}) {707};
\fill[cVerde] ({\vBx{2000}},{\vBy{43}}) circle (2.1pt);
\node[above right, font=\scriptsize, text=cVerde]
     at ({\vBx{2000}},{\vBy{43}}) {\SI{-57}{\percent} en 2000};
\fill[cVerde] ({\vBx{2026}},{\vBy{25}}) circle (2.1pt);
\node[above right, font=\scriptsize, text=cVerde]
     at ({\vBx{2026}},{\vBy{25}}) {\SI{-75}{\percent} en 2026};

\node[right, font=\scriptsize, text=cAviso, align=left]
     at ({\vBx{2004}},{\vBy{88}})
     {\SI{-2}{\percent} anual acumulativo:\\ mismo esfuerzo, otro eje};
\end{scope}

\end{tikzpicture}
\caption{La velocidad se detuvo en 1958; el esfuerzo tecnológico no. Panel A:
velocidades de crucero típicas a altitud de crucero. Entre el Boeing 707 de
1958 (\SI{897}{\kilo\meter\per\hour}) y el Airbus A350 de 2018
(\SI{903}{\kilo\meter\per\hour}) hay un 0,7 por ciento de diferencia en
sesenta años. El Concorde es el único punto que rompió la meseta y fue
retirado en 2003 \parencite{wikiconcorde}; el Overture, previsto para 2030, es
el intento de volver \parencite{wikioverture}. Panel B: índice de energía por
pasajero-kilómetro construido con la tasa de mejora de \SI{2}{\percent} anual
acumulativo que reporta la industria \parencite{iata2026eficiencia,peeters2005};
elaboración propia. Los dos paneles cubren el mismo periodo y cuentan la misma
historia vista desde dos figuras de mérito distintas.}
\label{fig:velocidad}
\end{figure}

Entre el Boeing 707 que entró en servicio en 1958 a unos
\SI{897}{\kilo\meter\per\hour} y el Airbus A350 que vuela hoy a
\SI{903}{\kilo\meter\per\hour} hay \textbf{un 0,7 por ciento de mejora en
sesenta y ocho años}. En los veintidós años anteriores, de 1936 a 1958, la
misma figura de mérito se había multiplicado por 2,7. La curva no es una S
incompleta: es una S que llegó a su asíntota y se quedó ahí.

\subsection{Por qué se detuvo, y por qué no fue un fracaso}

La razón física es conocida: cerca de Mach 1 la resistencia de onda se dispara,
y para atravesar esa barrera hace falta un ala completamente distinta, un motor
distinto y un consumo que se multiplica. La razón económica es la que decidió
el asunto. El Concorde cruzaba el Atlántico quemando prácticamente el mismo
combustible que un Boeing 747 para llevar cien pasajeros en lugar de
trescientos cincuenta \parencite{wikiconcorde}. Su consumo por asiento-milla
era entre dos veces y media y cinco veces peor que el de un subsónico
comparable. Voló veintisiete años y se retiró en 2003.

Y aquí está lo que a mí me parece la lección del módulo. \textbf{La industria
no dejó de innovar en 1958: cambió de figura de mérito.} El Panel B de la
Figura~\ref{fig:velocidad} muestra el otro eje, el del consumo de energía por
pasajero-kilómetro, mejorando a una tasa del orden del 2 por ciento anual
acumulativo de forma sostenida \parencite{iata2026eficiencia,peeters2005}. Un 2
por ciento anual durante sesenta y ocho años es un factor de cuatro: el avión
de hoy gasta la cuarta parte por pasajero y kilómetro que el de 1958.

Esto es exactamente lo que \textcite{deweck2022} describe como el
desplazamiento a lo largo de un frente de Pareto. Velocidad y consumo son dos
figuras de mérito en conflicto. El mercado eligió el punto de ese frente que
maximiza el asiento barato, no el asiento rápido, y lleva ahí desde los años
setenta. Mi queja, traducida al lenguaje del módulo, es que \emph{yo} preferiría
otro punto del frente. Es una preferencia legítima, pero es minoritaria: cada
vez que se le ha puesto precio, el pasajero medio ha elegido barato antes que
rápido.

\subsection{Qué pasa si de verdad se duplica la velocidad}

Aquí es donde mi respuesta espontánea se rompe, y el cálculo lo enseña mejor que
cualquier argumento. Tomo mi ruta larga, San José--Los Ángeles, y descompongo el
viaje entero de puerta a puerta, no solo el vuelo. La Figura~\ref{fig:puerta}
es el resultado.

\begin{figure}[t]
\centering
\begin{tikzpicture}

\node[anchor=south west, font=\small\bfseries] at (-3.3,3.72)
     {San José--Los Ángeles, \SI{4382}{\kilo\meter}: viaje completo de puerta
      a puerta};

% la meta que yo mismo me puse: la mitad de las diez horas de hoy
\draw[cVerde, dashed] ({5.0*1.02},-0.60) -- ({5.0*1.02},3.05);
\node[above, font=\scriptsize, text=cVerde, align=center]
     at ({5.0*1.02},3.02) {mi meta: la mitad};

\barrapp{2.55}{2.75}{5.5}{1.75}{\textbf{Hoy}\\Mach 0,85}{cSub}{10,0}
\barrapp{1.60}{2.75}{3.1}{1.75}{Overture\\Mach 1,7 (2030)}{cSup}{7,6}
\barrapp{0.65}{2.75}{2.8}{1.75}{Clase Concorde\\Mach 2,0}{cSup}{7,3}
\barrapp{-0.30}{1.40}{2.8}{0.85}{Mach 2,0 \emph{y} mitad
   de tiempo en tierra}{cSup}{5,1}

% leyenda
\fill[cTierra!45] (0,-1.20) rectangle (0.5,-0.95);
\node[right, font=\scriptsize] at (0.55,-1.07)
     {tiempo en tierra: traslados, facturación, seguridad, equipaje,
      inmigración};
\fill[cSub] (0,-1.60) rectangle (0.5,-1.35);
\fill[cSup] (0.5,-1.60) rectangle (1.0,-1.35);
\node[right, font=\scriptsize] at (1.05,-1.47)
     {tiempo dentro del avión, subsónico y supersónico};

\end{tikzpicture}
\caption{Duplicar la velocidad no divide el viaje entre dos. Descomposición del
viaje completo San José--Los Ángeles en tiempo en tierra
(\SI{2.75}{\hour} antes del despegue y \SI{1.75}{\hour} después del
aterrizaje, que es lo que a mí me cuesta de verdad) y tiempo de vuelo
(\SI{0.6}{\hour} fijas de rodaje, ascenso y descenso, más
\SI{0.7}{\hour} de penalización de aceleración y desaceleración en los casos
supersónicos, más la distancia dividida entre la velocidad de crucero).
Elaboración propia. Las \SI{4.5}{\hour} en tierra son idénticas en las tres
primeras filas: son la fracción que la velocidad no puede tocar.}
\label{fig:puerta}
\end{figure}

El viaje entero me lleva hoy unas diez horas de puerta a puerta, de las cuales
cinco y media están dentro del avión y cuatro y media están en tierra:
traslados, llegar dos horas antes, facturación, seguridad, desembarque, equipaje
e inmigración. Ese bloque de cuatro horas y media \textbf{es indiferente a la
velocidad del avión}.

De ahí sale el resultado incómodo:

\begin{itemize}
\item Con el Overture a Mach 1,7 el viaje pasa de 10,0 a 7,6 horas. Es un
  \textbf{24 por ciento} menos, no un 50.
\item Ni siquiera Mach 2,0 llega: 7,3 horas, un 27 por ciento menos. El
  Concorde no me habría dado lo que pedí.
\item Para que el viaje dure la mitad hace falta atacar \emph{a la vez} la
  velocidad y el tiempo en tierra. La cuarta fila de la
  Figura~\ref{fig:puerta}, con Mach 2,0 y el tiempo en tierra reducido a la
  mitad, es la única que llega a 5,1 horas.
\end{itemize}

Es la ley de Amdahl aplicada a un viaje \parencite{amdahl1967}: acelerar sin
límite una parte del proceso deja el total gobernado por la parte que no se
aceleró. Y la parte que no se aceleró aquí no es aeronáutica, es
administrativa: control de pasaportes, entrega de equipaje y el margen de
seguridad que yo mismo me doy al salir de casa. Buena parte de eso ya sabemos
cómo resolverlo, y es mucho más barato que un motor supersónico.

\subsection{La respuesta correcta a mi propia queja}

Al hacer el cálculo me di cuenta de que había declarado mal mi figura de
mérito. Yo dije \enquote{que el viaje dure la mitad}, pero lo que escribí a
continuación fue \enquote{me desespero luego de estar más de cinco horas en un
avión}. \textbf{No son la misma medida.} La primera es el tiempo puerta a
puerta; la segunda es el tiempo sentado en la cabina, y es esa la que de verdad
me molesta.

Y la buena noticia es que la segunda sí se resuelve con velocidad, y con mucho
menos de lo que yo pedía:

\begin{itemize}
\item Hoy, San José--Los Ángeles son \textbf{5,5 horas en cabina}: justo por
  encima de mi umbral.
\item A Mach 1,7 serían \textbf{3,1 horas}. Cruza el umbral con holgura.
\item San José--Houston, que hoy son 3,4 horas, bajaría a 2,0.
\end{itemize}

Es decir: \textbf{Mach 1,7 no me da la mitad del viaje, pero sí me quita
exactamente lo que me molesta.} Declarar bien la figura de mérito cambió el
requisito de \enquote{reducir el viaje un 50 por ciento}, que es probablemente
inalcanzable, a \enquote{mantener el tiempo en cabina por debajo de cuatro
horas}, que es un objetivo concreto y que un supersónico de Mach 1,7 cumple en
todas mis rutas.

Sobre si eso va a existir: el demostrador XB-1 de Boom rompió la barrera del
sonido el 28 de enero de 2025, el primer avión supersónico de desarrollo
privado que lo consigue \parencite{boom2025xb1}, y el Overture apunta a entrar
en servicio hacia 2030 \parencite{wikioverture}. En paralelo, el X-59 de la NASA
voló por primera vez el 28 de octubre de 2025 y superó la velocidad del sonido
el 5 de junio de 2026, a Mach 1,1 y \num{43400} pies
\parencite{nasa2025x59,nasa2026x59}. Su objetivo no es la velocidad sino el
ruido: reunir datos de respuesta de la población a un estampido atenuado para
que los reguladores puedan levantar la prohibición de volar en supersónico
sobre tierra firme, que es la restricción que dejó al Concorde encerrado en
rutas transatlánticas. \textbf{La barrera que hay que romper esta vez no es la
del sonido, es la regulatoria y la del consumo}, y el problema del consumo
sigue sin resolverse: sigue siendo el mismo frente de Pareto de 1976.

\section{Qué opino sobre la idea de la \enquote{vergüenza de volar}}
\label{sec:verguenza}

\textbf{A mí no me afecta en absoluto.} Cada viaje de trabajo que hago es
estrictamente necesario, y no siento culpa por hacerlo. Mantengo esa respuesta
después de haber hecho los números, pero los números me obligan a matizarla en
un punto, y prefiero decirlo yo antes de que me lo digan.

\subsection*{Qué es y qué efecto tuvo}

El término \emph{flygskam} nace en Suecia en 2017 y se convierte en un
movimiento social entre 2018 y 2019 \parencite{wikiflygskam}. Su efecto fue
real y medible: los vuelos internacionales en aeropuertos suecos cayeron un 4
por ciento en un año, los pasajeros domésticos un 9 por ciento y el operador
ferroviario estatal reportó un aumento del 21 por ciento en viajes en tren
\parencite{aiaa2020flygskam}. Es uno de los pocos casos documentados de un
movimiento de opinión que mueve una figura de mérito de demanda.

Ahora bien, ese efecto fue \textbf{local y transitorio}: se concentró en un
país con red ferroviaria densa, alternativa real y renta alta, y quedó
enterrado por la recuperación posterior a 2020. Como palanca global, la
vergüenza no ha movido la aguja.

\subsection*{El tamaño real del problema}

La aviación comercial representa alrededor del 2,5 por ciento de las emisiones
mundiales de \ce{CO2} \parencite{owid2026aviacion}. Si se contabiliza el
efecto climático completo, incluidas las estelas de condensación, los óxidos de
nitrógeno y el vapor de agua estratosférico, la cifra sube a cerca del
\textbf{3,5 por ciento del forzamiento radiativo antropogénico}, y ahí está el
dato que casi nadie menciona: \textbf{dos tercios de ese efecto no son
\ce{CO2}} \parencite{lee2021}. Solo las estelas de condensación pesan unas tres
veces más que todo el \ce{CO2} acumulado de la aviación.

Esto tiene una consecuencia que va en contra de la intuición del movimiento:
si dos tercios del daño vienen de las estelas, y las estelas dependen de la
altitud, la ruta y la humedad de la capa de aire que se atraviesa,
\textbf{desviar unos pocos vuelos unos pocos miles de pies puede valer más que
convencer a mucha gente de no volar}. Es un problema de ingeniería de
trayectorias, no de moral individual.

\begin{figure}[t]
\centering
\begin{tikzpicture}

% ---------- Panel A: mi huella frente a dos referencias ----------
\node[anchor=south west, font=\small\bfseries] at (-0.7,4.55)
     {Panel A. Mis seis tramos frente a dos referencias, en toneladas de
      \ce{CO2}e al año};

\foreach \v in {1,2,3,4} {
  \draw[cTierra!30] (0,{\v*0.86}) -- (13.9,{\v*0.86});
  \node[left, font=\scriptsize, text=cTierra] at (-0.08,{\v*0.86}) {\v};
}
\draw[cTierra!70] (0,0) -- (13.9,0);

\barrav{2.0}{1.385}{cTierra}{media anual \emph{total} de un costarricense}{1,61}
\barrav{5.4}{1.938}{cSup}{mis 6 tramos, todos a Houston}{2,3}
\barrav{8.5}{2.595}{cSup}{mis 6 tramos, todos a Los Ángeles}{3,1}
\barrav{11.9}{4.042}{cTierra}{media anual mundial por persona}{4,7}

\draw[cAviso, dashed] (0,1.385) -- (13.9,1.385);
\node[right, font=\scriptsize\itshape, text=cAviso]
     at (12.9,1.60) {};

% ---------- Panel B: la aritmética que la vergüenza no gana ----------
\begin{scope}[yshift=-6.1cm]
\node[anchor=south west, font=\small\bfseries] at (-0.7,3.95)
     {Panel B. La aritmética de fondo: eficiencia contra tráfico
      (índice, 2019 = 100)};

\foreach \v in {50,100,150,200} {
  \draw[cTierra!30] (0,{\eAy{\v}}) -- (13.9,{\eAy{\v}});
  \node[left, font=\scriptsize, text=cTierra] at (-0.08,{\eAy{\v}}) {\v};
}
\draw[cTierra!70] (0,0) -- (13.9,0);
\draw[cTierra!70] (0,0) -- (0,3.75);
\foreach \a in {2010,2015,2020,2025,2030,2035,2040} {
  \draw[cTierra!70] ({\eAx{\a}},0) -- ({\eAx{\a}},-0.09);
  \node[below, font=\scriptsize, text=cTierra] at ({\eAx{\a}},-0.09) {\a};
}

\draw[cAviso, very thick, samples=60, domain=2010:2040, variable=\x]
  plot ({\eAx{\x}},{(100*pow(1.036,\x-2019))*0.0160});
\node[right, font=\scriptsize, text=cAviso, align=left]
     at ({\eAx{2040}+0.12},{\eAy{205}}) {tráfico\\$+3{,}6$ \% anual};

\draw[cRef, very thick, samples=60, domain=2010:2040, variable=\x]
  plot ({\eAx{\x}},{(100*pow(1.0153,\x-2019))*0.0160});
\node[right, font=\scriptsize, text=cRef, align=left]
     at ({\eAx{2040}+0.12},{\eAy{137}}) {emisiones\\$+1{,}5$ \% anual};

\draw[cVerde, very thick, samples=60, domain=2010:2040, variable=\x]
  plot ({\eAx{\x}},{(100*pow(0.98,\x-2019))*0.0160});
\node[right, font=\scriptsize, text=cVerde, align=left]
     at ({\eAx{2040}+0.12},{\eAy{65}}) {eficiencia\\$-2$ \% anual};

\draw[cTierra!80, dashed] ({\eAx{2019}},0) -- ({\eAx{2019}},{\eAy{112}});
\node[above, font=\scriptsize, text=cTierra] at ({\eAx{2019}},{\eAy{112}})
     {2019};
\end{scope}

\end{tikzpicture}
\caption{Los dos números que hacen falta para opinar sobre la vergüenza de
volar. Panel A: seis tramos de mis rutas habituales emiten entre 2,3 y 3,1
toneladas de \ce{CO2} equivalente, calculadas con los factores de
\textcite{defra2025} sobre las distancias del Cuadro~\ref{cuadro:rutas}. La
línea roja es la emisión media anual \emph{total} de un costarricense, 1,61
toneladas \parencite{owid2026costarica}. Panel B: la eficiencia por
pasajero-kilómetro mejora en torno al 2 por ciento anual
\parencite{iata2026eficiencia} y el tráfico crece más deprisa; el producto de
las dos series es la línea de emisiones, que sigue subiendo. Elaboración
propia.}
\label{fig:huella}
\end{figure}

\subsection*{Dónde tengo que matizar mi respuesta}

El Panel A de la Figura~\ref{fig:huella} es el dato que me obliga a matizar.
Mis seis tramos emiten entre 2,3 y 3,1 toneladas de \ce{CO2} equivalente. La
emisión media anual \emph{total} de un costarricense, contando todo lo que
hace en el año, es de 1,61 toneladas \parencite{owid2026costarica}. Es decir:
\textbf{seis vuelos de trabajo pesan entre 1,4 y 1,9 veces lo que emite un
costarricense medio en un año entero.}

Eso no me hace cambiar de decisión, y sostengo que los viajes eran necesarios.
Pero sí me obliga a abandonar un argumento que había usado sin pensarlo: el de
que mi contribución es despreciable. No lo es. Lo correcto es decir que es
\emph{necesaria}, que es otra cosa y es un argumento más honesto.

\subsection*{Por qué sigo pensando que la vergüenza no es el instrumento}

El Panel B de la Figura~\ref{fig:huella} contiene la razón de fondo. La
eficiencia por pasajero-kilómetro mejora en torno al 2 por ciento anual
\parencite{iata2026eficiencia}, y el tráfico crece por encima del 3,5 por
ciento anual. El producto de las dos series es una línea de emisiones que sube
alrededor de un 1,5 por ciento al año. \textbf{La mejora tecnológica va
ganando la carrera por pasajero y perdiéndola en total}, que es la definición
de un problema que no se arregla por el lado de la eficiencia marginal.

Contra esa aritmética, el mejor resultado documentado de la vergüenza de volar
es un 4 por ciento en un país durante un año. La conclusión no es que el
movimiento sea inútil: es que \textbf{no está apuntando a la figura de mérito
que gobierna el resultado}. Las palancas con la magnitud correcta son otras y
todas son tecnológicas o regulatorias: combustible de aviación sostenible,
evitación de estelas por gestión de trayectoria, y un precio del carbono que
entre en el billete. La vergüenza sirve, en todo caso, como señal de demanda
que mueve capital hacia esas tres. Ese es su papel útil, y no es pequeño, pero
es indirecto.

\section{Los tres mejores aeropuertos del mundo}

Aquí tengo que empezar reconociendo el límite: \textbf{no conozco muchos
aeropuertos.} Mi respuesta espontánea fueron dos, Atlanta y Panamá, y al
mirarla con el instrumental del módulo resulta que estaba usando \emph{dos
figuras de mérito distintas sin darme cuenta}, y que ninguna de las dos es la
que usan los rankings publicados.

\begin{table}[t]
\small
\caption{Cuatro figuras de mérito, cuatro podios distintos. Ninguna es
incorrecta: miden funciones distintas del mismo sistema.}
\label{cuadro:aeropuertos}
\begin{tabularx}{\textwidth}{@{}p{3.5cm}p{3.6cm}XX@{}}
\toprule
\textbf{Figura de mérito} & \textbf{Qué mide} & \textbf{Podio} &
\textbf{Fuente} \\
\midrule
Tránsito de pasajeros
& Volumen del nodo
& Atlanta (106,3 M), Dubái (95,2 M), Tokio Haneda (91,7 M)
& \textcite{aci2026transito} \\
Movimientos de aeronaves
& Densidad de operación
& Chicago O'Hare (\num{860015}), Atlanta (\num{807625}), Dallas
& \textcite{ohare2026} \\
Experiencia del pasajero
& Calidad del paso por la terminal
& Singapur Changi, Seúl Incheon, Tokio Haneda
& \textcite{skytrax2026} \\
Conectividad de red
& Destinos alcanzables sin segunda escala
& Panamá Tocumen: 88 destinos en 32 países
& \textcite{copa2025} \\
\bottomrule
\end{tabularx}
\end{table}

\subsection*{Qué estaba midiendo yo}

\textbf{Con Atlanta yo estaba midiendo tránsito y densidad de operación.} Dije
que me impresiona \enquote{la cantidad de vuelos que veo salir desde las
terminales}, y la impresión es correcta en magnitud: Atlanta movió 106,3
millones de pasajeros en 2025 y sigue siendo el aeropuerto con más tránsito del
mundo, título que ha tenido en veintisiete de los últimos veintiocho años
\parencite{aci2026transito}. Con una precisión que a mí me sorprendió: en
número de \emph{despegues y aterrizajes}, que es lo que yo estaba mirando de
verdad desde la terminal, Atlanta ya no es el primero. En 2025 lo adelantó
Chicago O'Hare con \num{860015} operaciones frente a \num{807625}, un
despegue o aterrizaje cada treinta y siete segundos durante todo el año
\parencite{ohare2026}. Mi intuición apuntaba al aeropuerto correcto para la
figura de mérito equivocada por un puesto.

\textbf{Con Panamá yo estaba midiendo otra cosa completamente distinta.} Lo
dije sin rodeos: \enquote{obviamente no es de los mejores del mundo, pero me
ayuda a conectarme con muchos destinos donde no tengo vuelos directos}. Eso no
es calidad de terminal, es \textbf{conectividad de red}, y en esa figura de
mérito Tocumen es excelente: el Hub de las Américas conecta 88 destinos en 32
países con más de 375 vuelos diarios \parencite{copa2025}, y movió más de 21
millones de pasajeros en 2025 \parencite{tocumen2025}. Para alguien que sale
de San José, Tocumen no compite con Changi: compite con no poder llegar.

\subsection*{Mi tercero, y por qué lo declaro distinto}

Como tercero pongo \textbf{Singapur Changi}, y lo marco explícitamente: no lo
conozco. Lo incluyo porque es el primero en la única figura de mérito que a mí
me falta por experiencia propia, la de la calidad del paso por la terminal, y
lo es por decimocuarta vez en los premios Skytrax de 2026
\parencite{skytrax2026}. Es una respuesta de segunda mano y prefiero decirlo
que fingir un criterio que no tengo.

La lección que me llevo del Cuadro~\ref{cuadro:aeropuertos} es la del módulo,
en pequeño: \textbf{cuatro figuras de mérito legítimas producen cuatro podios
sin un solo nombre en común entre el primero y el cuarto.} La pregunta
\enquote{cuál es el mejor aeropuerto} no tiene respuesta hasta que alguien diga
qué está midiendo.

\section{Los tres aeropuertos que más frecuento}

Son estos tres, y el orden importa poco porque los tres son el mismo viaje:

\begin{enumerate}
\item \textbf{Juan Santamaría (SJO)}, San José. De ahí salgo siempre.
\item \textbf{George Bush Intercontinental (IAH)}, Houston. Vuelo directo
  desde Costa Rica.
\item \textbf{Los Ángeles (LAX)}. Vuelo directo desde Costa Rica.
\end{enumerate}

Al escribirlo veo que esta lista \textbf{no expresa ninguna preferencia mía}.
Es la topología de la red vista desde San José. Juan Santamaría está porque
vivo aquí, y Houston y Los Ángeles están \emph{porque tienen vuelo directo}. La
razón que yo mismo di al nombrarlos es la razón entera.

\begin{figure}[t]
\centering
\begin{tikzpicture}[
  aer/.style={draw=cRef, fill=white, rounded corners=2pt, inner sep=3.5pt,
              font=\scriptsize, align=center},
  yo/.style={draw=cSup, very thick, fill=cSup!8, rounded corners=2pt,
             inner sep=4pt, font=\scriptsize\bfseries, align=center}]

\node[yo] (sjo) at (0,0) {Juan Santamaría\\(SJO)\\
  \normalfont\num{6403220} pas. en 2025};

\node[aer] (iah) at (7.1,2.05) {Houston\\(IAH)};
\node[aer] (lax) at (7.1,0.0)  {Los Ángeles\\(LAX)};
\node[aer, draw=cSub, fill=cSub!8] (pty) at (7.1,-2.05)
     {Panamá Tocumen\\(PTY)};

\draw[cSup, very thick, -{Stealth[length=5pt]}] (sjo) -- (iah)
  node[pos=0.46, above, sloped, font=\scriptsize, text=cSup]
  {directo, \SI{2505}{\kilo\meter}, 3,4 h};
\draw[cSup, very thick, -{Stealth[length=5pt]}] (sjo) -- (lax)
  node[pos=0.5, above, font=\scriptsize, text=cSup]
  {directo, \SI{4382}{\kilo\meter}, 5,5 h};
\draw[cSub, very thick, -{Stealth[length=5pt]}] (sjo) -- (pty)
  node[pos=0.46, below, sloped, font=\scriptsize, text=cSub]
  {directo, \SI{539}{\kilo\meter}, 1,2 h};

\foreach \a in {-38,-25,-12,1,14,27,40} {
  \draw[cSub!55, -{Stealth[length=3.5pt]}]
    (pty.east) -- ++(\a:1.15);
}
\node[right, font=\scriptsize, text=cSub, align=left] at (9.75,-2.05)
     {88 destinos\\en 32 países\\\emph{sin vuelo directo}};

\node[font=\scriptsize\itshape, text=cTierra, align=center]
     at (5.0,-3.55)
     {Los dos destinos que más frecuento son los dos que tienen vuelo directo.
      No es una preferencia, es la red.};

\end{tikzpicture}
\caption{Mi red de vuelos, que explica por sí sola cuáles son mis tres
aeropuertos. Houston y Los Ángeles aparecen porque son directos desde San
José; Panamá aparece porque convierte un destino inalcanzable en uno de dos
tramos. Cifras de pasajeros de \textcite{aeris2026} y de destinos de
\textcite{copa2025}; distancias y tiempos, cálculo propio (Cuadro~\ref{cuadro:rutas}).}
\label{fig:red}
\end{figure}

Y esto conecta con la pregunta anterior de una forma que no había visto:
\textbf{los aeropuertos que más frecuento no son los que mejor me parecen, y
los que mejor me parecen no son los mejores según ningún ranking.} La figura de
mérito que gobierna mi comportamiento real es la conectividad directa desde San
José, no la calidad de la terminal ni el tamaño del nodo. Juan Santamaría, por
cierto, cerró 2025 con \num{6403220} pasajeros, el mejor año de su historia y
un 3 por ciento más que el anterior \parencite{aeris2026}: es un aeropuerto
pequeño en la escala del Cuadro~\ref{cuadro:aeropuertos} y es, para mí, el
único que no puedo sustituir.

\section{Modos de transporte más lentos con menor huella de carbono}

Mi respuesta fue esta: \textbf{es algo de lo que no debería preocuparme; hay
que dejar que la tecnología avance por sí sola.} La mantengo en su parte
práctica y la corrijo en su parte causal.

\subsection*{Dónde tengo razón: en Costa Rica la pregunta está vacía}

Para mis rutas, el frente de Pareto entre tiempo y emisiones \textbf{tiene un
solo punto}. No hay tren a Houston. No hay tren a Los Ángeles. Ni siquiera hay
tren a Panamá. El corredor ferroviario que haría posible la sustitución no
existe en Centroamérica y no está en construcción. Cuando una pregunta de
elección tecnológica se plantea sobre un conjunto de alternativas de tamaño
uno, no es una elección: es una descripción.

Donde la pregunta sí tiene contenido es en Europa continental o en el corredor
Tokio--Osaka, con distancias por debajo de unos \SI{800}{\kilo\meter}, un tren
de alta velocidad que compite en tiempo puerta a puerta y una emisión por
pasajero-kilómetro que puede ser veinte veces menor \parencite{defra2025}. Ahí
la sustitución no exige sacrificio: el tren gana en las dos figuras de mérito a
la vez. Ese es el caso en el que la vergüenza de volar tuvo efecto medible, y
no es casualidad que fuera Suecia.

Dicho esto, contesto la pregunta tal como está formulada: \textbf{sí, usaría un
modo más lento si la diferencia de tiempo puerta a puerta fuera de una o dos
horas} y la reducción de emisiones fuera del orden de un factor diez. Si la
diferencia fuera de un día, no. Es la misma tasa de canje que aplico en
cualquier otra decisión.

\subsection*{Dónde me corrijo: la tecnología no avanza sola}

La parte de mi respuesta que no sostengo después de hacer este trabajo es la de
\enquote{dejar que la tecnología avance por sí sola}. El Panel B de la
Figura~\ref{fig:huella} muestra por qué. La eficiencia lleva sesenta años
mejorando al 2 por ciento anual, que es una tasa notable y sostenida, y aun así
\emph{pierde} contra el crecimiento del tráfico. Una tendencia que va
ganando la carrera equivocada no se corrige sola por el hecho de continuar.

Y sobre todo: esa tasa del 2 por ciento no es un fenómeno natural. Es el
resultado acumulado de programas de investigación, de normas de certificación,
de regulación de ruido y de emisiones, y de un precio del combustible que ha
hecho que el consumo sea la variable económica dominante de una aerolínea.
\textbf{La tecnología avanzó porque alguien puso el dinero y la norma en la
figura de mérito correcta.} La velocidad de crucero de la
Figura~\ref{fig:velocidad} es la demostración por contraste: es una figura de
mérito que también podría haber seguido mejorando y lleva sesenta y ocho años
parada, precisamente porque nadie puso ahí ni el dinero ni la norma.

Así que mi respuesta corregida es: no me corresponde a mí resolverlo dejando de
volar seis veces al año, en eso tenía razón. Pero decir que la tecnología
avanza sola es exactamente el tipo de afirmación que este módulo desmiente.
Avanza donde se la empuja, y la elección de dónde empujar es una decisión de
hoja de ruta.

\section{Lo que dicen juntas las seis respuestas}

Al ponerlas en fila aparece un patrón que no esperaba encontrar en un foro
sobre experiencias personales.

\begin{enumerate}
\item \textbf{En tres de las seis preguntas yo había declarado mal la figura de
  mérito.} Pedí velocidad cuando quería menos tiempo sentado; llamé
  \enquote{mejores} a dos aeropuertos midiendo dos cosas distintas; y llamé
  preferencia a lo que era topología de red. Declarar bien la medida cambió la
  respuesta en los tres casos.
\item \textbf{La figura de mérito que se satura no detiene el progreso: lo
  desvía.} La velocidad se paró en 1958 y el esfuerzo se fue al consumo, que
  desde entonces ha mejorado un factor de cuatro. Quien mire solo el eje de la
  velocidad concluirá que la aviación lleva sesenta años sin innovar, y estará
  midiendo el eje equivocado.
\item \textbf{Acelerar una parte no acelera el todo.} Mi mejora soñada, el doble
  de velocidad, produce un 24 por ciento de reducción del viaje, no un 50, y
  esto se sabe sin conocer nada de aerodinámica: basta con descomponer el
  proceso y ver qué fracción es sensible a la palanca.
\item \textbf{El instrumento tiene que apuntar a la figura de mérito que
  gobierna el resultado.} La vergüenza de volar actúa sobre la demanda con un
  efecto documentado del 4 por ciento en un país; la aritmética que hay que
  vencer es del 1,5 por ciento anual compuesto y global, y dos tercios del daño
  ni siquiera son \ce{CO2}.
\end{enumerate}

Si tuviera que quedarme con una sola frase: \textbf{casi todas mis quejas de
pasajero eran correctas en el síntoma y erróneas en la medida}, y el módulo
sirve justamente para eso, para pasar de lo primero a lo segundo.

\section{Limitaciones}

\begin{enumerate}
\item \textbf{No llevo registro de mis itinerarios.} Las seis veces que volé
  este año las he modelado como seis tramos sobre mis rutas habituales, y por
  eso doy un rango (2,3 a 3,1 toneladas) y no una cifra. Si los seis fueron
  tres viajes de ida y vuelta, el rango es el mismo; si fueron seis viajes de
  ida y vuelta, hay que duplicarlo.
\item \textbf{Las distancias son de círculo máximo}, no rutas reales. Las rutas
  reales son entre un 2 y un 6 por ciento más largas por el encaminamiento del
  tráfico aéreo, de modo que tiempos y emisiones están ligeramente
  subestimados.
\item \textbf{Los tiempos en tierra de la Figura~\ref{fig:puerta} son mis
  hábitos, no un promedio.} Dos horas de antelación en San José y una hora de
  salida en Los Ángeles son mis números; quien facture solo con equipaje de
  mano y viaje con precheque tendrá un bloque fijo menor y, por tanto, obtendrá
  más beneficio de la velocidad que yo.
\item \textbf{El Panel B de la Figura~\ref{fig:velocidad} es una curva
  construida}, no una serie observada: aplica la tasa de mejora que reporta la
  industria a un índice anclado en 1958. Sirve para ver la forma y el orden de
  magnitud, no para leer un valor concreto de un año concreto.
\item \textbf{Los factores de emisión son promedios de flota} en clase turista
  y no distinguen modelo de avión, factor de ocupación real ni efectos no
  \ce{CO2}. Si se incluyeran esos últimos con el multiplicador de
  \textcite{lee2021}, mis cifras del Panel A habría que multiplicarlas por un
  factor cercano a tres.
\item \textbf{El Overture no existe todavía.} Sus \SI{1806}{\kilo\meter\per\hour}
  son una especificación de proyecto, no un dato medido, y la fecha de 2030 es
  un objetivo del fabricante. El programa aún no tiene motor certificado.
\item \textbf{Los rankings de aeropuertos no son comparables entre sí.} El de
  Skytrax es una encuesta de satisfacción, los de ACI son recuentos
  operativos y el de conectividad lo he construido yo con datos de una sola
  aerolínea. Van juntos en el Cuadro~\ref{cuadro:aeropuertos} para mostrar que
  discrepan, no para promediarlos.
\item \textbf{No he conocido Changi.} Lo incluyo por el ranking y lo declaro
  como tal.
\end{enumerate}

\printbibliography

\end{document}
